% GENERATED FILE. Edit canonical lessons, then run npm run generate:books.
% canonical-source-hash: fnv1a64:ca80a83c22483785
% canonical-lessons: PT-C427-caucao, PT-C427-desmarcar, PT-C427-cancelamento, PT-C427-antecipar, PT-C427-remarcar

\chapter{Reservar e desmarcar}
\label{ch:pt-reservar-e-desmarcar}
\hlchaptermodality{427}

\begin{chapteropening}
\textbf{By the end of this chapter:} \emph{I can handle a booking: the deposit, calling it off, the cancellation, bringing dates forward and rebooking.}

Complete the last lesson of chapter 427: I can handle a booking: the deposit, calling it off, the cancellation, bringing dates forward and rebooking.
\end{chapteropening}

\section[a caução]{a caução — the security deposit; bail}

\label{lesson:PT-C427-caucao}



\begin{quote}
Before the new one: say the Portuguese for the tent peg, then the Portuguese for the mosquito net.
\end{quote}

\subsection*{You'll want to know: a caução}
\textbf{a caução} — \textquotedblleft{}the security deposit; bail\textquotedblright{}.

\begin{quote}
Pagámos uma caução de cem euros pela casa.
\end{quote}

\begin{quote}
\emph{We paid a hundred-euro deposit for the house.}
\end{quote}

You get \emph{a caução} back at the end if nothing is broken.

\begin{cousinweb}
\emph{caução} is a learned borrowing from Latin \emph{cautiō} 'caution, guarantee' (EtymDB), built on \emph{cavēre} 'to take care'.
\end{cousinweb}

\begin{grammarlens}[title={What you've built}]
One more word for telling, explaining and describing at B1.
\end{grammarlens}

\subsection*{Your turn}
\begin{itemize}
\raggedright
  \item \emph{Say it:} \emph{a caução}
  \item \emph{Say it:} \emph{a caução}, once more
  \item \emph{Say it:} a sentence of your own with \emph{a caução}
  \item \emph{Recall:} say the Portuguese for the tent peg, then the Portuguese for the mosquito net, then say \emph{a caução} again
\end{itemize}

\begin{tcolorbox}[breakable,colback=teal!4,colframe=teal!35!black,title={Before you move on}]
What does a caução mean? (\textquotedblleft{}The security deposit; bail\textquotedblright{}.) Say it once more.
\end{tcolorbox}
\section[desmarcar]{desmarcar — to call off, to cancel (an appointment or booking)}

\label{lesson:PT-C427-desmarcar}



\begin{quote}
Before the new one: say the Portuguese for the mosquito net, then the Portuguese for the security deposit.
\end{quote}

\subsection*{You'll want to know: desmarcar}
\textbf{desmarcar} — \textquotedblleft{}to call off, to cancel (an appointment or booking)\textquotedblright{}.

\begin{quote}
Desmarquei o quarto porque fiquei doente.
\end{quote}

\begin{quote}
\emph{I cancelled the room because I fell ill.}
\end{quote}

\emph{Marcar} — to book; \emph{desmarcar} — to call it off.

\begin{cousinweb}
\emph{desmarcar} is built from \emph{des-} and \emph{marcar} 'to mark, to book'. \emph{Marcar} is usually linked with \emph{marco} 'boundary stone', from Germanic \emph{markō} 'boundary' (EtymDB).
\end{cousinweb}

\begin{grammarlens}[title={What you've built}]
One more word for telling, explaining and describing at B1.
\end{grammarlens}

\subsection*{Your turn}
\begin{itemize}
\raggedright
  \item \emph{Say it:} \emph{desmarcar}
  \item \emph{Say it:} \emph{desmarcar}, once more
  \item \emph{Say it:} a sentence of your own with \emph{desmarcar}
  \item \emph{Recall:} say the Portuguese for the mosquito net, then the Portuguese for the security deposit, then say \emph{desmarcar} again
\end{itemize}

\begin{tcolorbox}[breakable,colback=teal!4,colframe=teal!35!black,title={Before you move on}]
What does desmarcar mean? (\textquotedblleft{}To call off, to cancel (an appointment or booking)\textquotedblright{}.) Say it once more.
\end{tcolorbox}
\section[o cancelamento]{o cancelamento — the cancellation}

\label{lesson:PT-C427-cancelamento}



\begin{quote}
Before the new one: say the Portuguese for the security deposit, then the Portuguese for to call off.
\end{quote}

\subsection*{You'll want to know: o cancelamento}
\textbf{o cancelamento} — \textquotedblleft{}the cancellation\textquotedblright{}.

\begin{quote}
O cancelamento é gratuito até dois dias antes.
\end{quote}

\begin{quote}
\emph{Cancellation is free up to two days before.}
\end{quote}

The verb is \emph{cancelar}.

\begin{cousinweb}
\emph{cancelamento} is built on \emph{cancelar} 'to cancel', a learned borrowing from Latin \emph{cancellāre} 'to cross out with lattice lines', on \emph{cancellī} 'lattice, railings' (Lewis \& Short).
\end{cousinweb}

\begin{grammarlens}[title={What you've built}]
One more word for telling, explaining and describing at B1.
\end{grammarlens}

\subsection*{Your turn}
\begin{itemize}
\raggedright
  \item \emph{Say it:} \emph{o cancelamento}
  \item \emph{Say it:} \emph{o cancelamento}, once more
  \item \emph{Say it:} a sentence of your own with \emph{o cancelamento}
  \item \emph{Recall:} say the Portuguese for the security deposit, then the Portuguese for to call off, then say \emph{o cancelamento} again
\end{itemize}

\begin{tcolorbox}[breakable,colback=teal!4,colframe=teal!35!black,title={Before you move on}]
What does o cancelamento mean? (\textquotedblleft{}The cancellation\textquotedblright{}.) Say it once more.
\end{tcolorbox}
\section[antecipar]{antecipar — to bring forward; to anticipate}

\label{lesson:PT-C427-antecipar}



\begin{quote}
Before the new one: say the Portuguese for to call off, then the Portuguese for the cancellation.
\end{quote}

\subsection*{You'll want to know: antecipar}
\textbf{antecipar} — \textquotedblleft{}to bring forward; to anticipate\textquotedblright{}.

\begin{quote}
Antecipámos o regresso por causa do tempo.
\end{quote}

\begin{quote}
\emph{We brought our return forward because of the weather.}
\end{quote}

\emph{Com antecedência} — in advance.

\begin{cousinweb}
\emph{antecipar} is a learned borrowing from Latin \emph{anticipāre} 'to take beforehand' (EtymDB), from \emph{ante} 'before' and \emph{capere} 'to take'.
\end{cousinweb}

\begin{grammarlens}[title={What you've built}]
One more word for telling, explaining and describing at B1.
\end{grammarlens}

\subsection*{Your turn}
\begin{itemize}
\raggedright
  \item \emph{Say it:} \emph{antecipar}
  \item \emph{Say it:} \emph{antecipar}, once more
  \item \emph{Say it:} a sentence of your own with \emph{antecipar}
  \item \emph{Recall:} say the Portuguese for to call off, then the Portuguese for the cancellation, then say \emph{antecipar} again
\end{itemize}

\begin{tcolorbox}[breakable,colback=teal!4,colframe=teal!35!black,title={Before you move on}]
What does antecipar mean? (\textquotedblleft{}To bring forward; to anticipate\textquotedblright{}.) Say it once more.
\end{tcolorbox}
\section[remarcar]{remarcar — to rebook, to reschedule}

\label{lesson:PT-C427-remarcar}



\begin{quote}
Before the new one: say the Portuguese for the cancellation, then the Portuguese for to bring forward.
\end{quote}

\subsection*{You'll want to know: remarcar}
\textbf{remarcar} — \textquotedblleft{}to rebook, to reschedule\textquotedblright{}.

\begin{quote}
A agência remarcou o voo para sexta.
\end{quote}

\begin{quote}
\emph{The agency rebooked the flight for Friday.}
\end{quote}

Also to mark a price again.

\begin{cousinweb}
\emph{remarcar} is built from \emph{re-} 'again' and \emph{marcar} 'to mark, to book', which is usually linked with \emph{marco} 'boundary stone', from Germanic \emph{markō} 'boundary' (EtymDB).
\end{cousinweb}

\begin{grammarlens}[title={What you've built}]
One more word for telling, explaining and describing at B1.
\end{grammarlens}

\subsection*{Your turn}
\begin{itemize}
\raggedright
  \item \emph{Say it:} \emph{remarcar}
  \item \emph{Say it:} \emph{remarcar}, once more
  \item \emph{Say it:} a sentence of your own with \emph{remarcar}
  \item \emph{Recall:} say the Portuguese for the cancellation, then the Portuguese for to bring forward, then say \emph{remarcar} again
\end{itemize}

\begin{tcolorbox}[breakable,colback=teal!4,colframe=teal!35!black,title={Before you move on}]
What does remarcar mean? (\textquotedblleft{}To rebook, to reschedule\textquotedblright{}.) Say it once more.
\end{tcolorbox}
